% generator_I02.tex -- Prop I.2 (Prob): from a given point, draw a line equal to a given line.
% Figure and proof geometry from jemmybutton byrne-en-latex.tex (CC-BY-SA), wording from 1847 original.
% Build: lualatex generator_I02.tex  ->  generator_I02.pdf
\documentclass[12pt]{article}
\usepackage{geometry}
\geometry{a5paper, margin=1.5cm}
\usepackage[lining]{ebgaramond}
\usepackage{amssymb}
\usepackage{byrne}
\def\mpPre{u := 1cm; textLabels := false;}

\begin{document}

% FIGURE: point A (given), line BC (given), equilateral triangle ABD, circles at D and B.
\defineNewPicture{%
    pair A, B, C, D, E, F;%
    path P[];%
    numeric r[];%
    A := (0, 0);%
    B := (-3/5u, -3/5u);%
    C := (-2u, -1/3u);%
    r1 := abs(A-B);%
    D := (fullcircle scaled 2r1 shifted A) intersectionpoint (fullcircle scaled 2r1 shifted B);%
    r2 := abs(B-C);%
    r3 := r1 + r2;%
    P1 := fullcircle scaled 2r2 shifted B;%
    P2 := fullcircle scaled 2r3 shifted D;%
    E := (D -- 10[D, B]) intersectionpoint P1;%
    F := (D -- 10[D, A]) intersectionpoint P2;%
    % Layer-1 assertions: AF = BC (the construction goal)
    if abs(abs(A-F) - abs(B-C)) > 0.01: errmessage "ASSERT: AF != BC"; fi;%
    % BD = DA (equilateral triangle sides)
    if abs(abs(B-D) - abs(D-A)) > 0.01: errmessage "ASSERT: BD != DA"; fi;%
    byLineDefine(A, B, byblack, DASHED_LINE, REGULAR_WIDTH);%
    byLineDefine(B, C, byblack, SOLID_LINE, REGULAR_WIDTH);%
    byLineDefine(B, D, byred,   SOLID_LINE, REGULAR_WIDTH);%
    byLineDefine(D, A, byred,   SOLID_LINE, REGULAR_WIDTH);%
    byLineDefine(B, E, byyellow,SOLID_LINE, REGULAR_WIDTH);%
    byLineDefine(A, F, byblue,  SOLID_LINE, REGULAR_WIDTH);%
    draw byNamedLineSeq(0)(AB,BC);%
    draw byNamedLineSeq(0)(BE,BD,DA,AF);%
    draw byCircle.A(D, E, byred,  0, 0,  1/2);%
    draw byCircle.B(B, C, byblue, 0, 0, -1/2);%
    draw byLabelsOnPolygon(E, D, A, F)(OMIT_FIRST_LABEL+OMIT_LAST_LABEL, 1);%
    draw byLabelsOnPolygon(E, B, C)(OMIT_FIRST_LABEL+OMIT_LAST_LABEL, 1);%
    draw byLabelsOnCircle(C)(B);%
    draw byLabelsOnCircle(E, F)(A);%
}

% STATEMENT
\begin{center}
\textbf{Book~I.\enspace Prop.\ II. Prob.}
\end{center}

\vspace{0.3cm}
\noindent
\begin{minipage}[t]{0.50\linewidth}
\itshape From a given point
(\drawFromCurrentPicture[middle][pointA]{%
    startGlobalRotation(-lineAngle.DA);%
    draw byNamedPointLines(A,"AB");%
    stopGlobalRotation;%
}),
to draw a straight line equal to a given straight line (\drawUnitLine{BC}).
\end{minipage}%
\hfill
\begin{minipage}[t]{0.46\linewidth}
\centering
\drawCurrentPicture
\end{minipage}

\vspace{0.5cm}

% PROOF
\begin{center}
Draw \drawUnitLine{AB}~(post.~I), describe%
\offsetPicture{10pt}{0pt}{%
\drawFromCurrentPicture[bottom]{%
    startAutoLabeling;%
    startTempScale(scaleFactor*3);%
    startGlobalRotation(180-lineAngle.AB);%
    draw byNamedLineSeq(0)(AB,BD,DA);%
    stopGlobalRotation;%
    stopTempScale;%
    stopAutoLabeling;%
}}%
(pr.~I.1),\\[3pt]
produce \drawUnitLine{BD}~(post.~II),\\[3pt]
describe \offsetPicture{10pt}{0pt}{\drawFromCurrentPicture{%
    draw byNamedLine(BC);%
    draw byNamedCircle(B);%
    draw byLabelLineEnd(B, C, 0);%
    draw byLabelLineEnd(C, B, 0);%
}} and \offsetPicture{10pt}{0pt}{\drawFromCurrentPicture{%
    draw byNamedLineSeq(0)(BD, BE);%
    draw byNamedCircle(A);%
    draw byLabelLineEnd(D, E, 0);%
    draw byLabelLineEnd(E, D, 1);%
}}~(post.~III);\\[3pt]
produce \drawUnitLine{DA}~(post.~II),\\[3pt]
then \drawUnitLine{AF} is the line required.\\[8pt]
For $\drawUnitLine{BE,BD} = \drawUnitLine{DA,AF}$~(def.~15),\\
and $\drawUnitLine{BD} = \drawUnitLine{DA}$~(const.),\\
$\therefore \drawUnitLine{BE} = \drawUnitLine{AF}$~(post.~III),\\
but~(def.~15) $\drawUnitLine{BC} = \drawUnitLine{BE} = \drawUnitLine{AF}$;\\[6pt]
$\therefore \drawUnitLine{AF}$ drawn from the given point~(\pointA),
is equal to the given line~\drawUnitLine{BC}.
\end{center}

\begin{flushright}Q.\ E.\ D.\end{flushright}

\end{document}
